% ============================================================
% Phase-Filtered Ramanujan Sums and the Spectral Gate
% ============================================================

\documentclass[12pt]{amsart}

\usepackage{amsmath,amssymb,amsthm}
\usepackage{booktabs}
\usepackage{microtype}
\usepackage{url}

\newtheorem{theorem}{Theorem}
\newtheorem{lemma}[theorem]{Lemma}
\newtheorem{proposition}[theorem]{Proposition}
\newtheorem{corollary}[theorem]{Corollary}
\newtheorem{conjecture}[theorem]{Conjecture}
\theoremstyle{definition}
\newtheorem{definition}[theorem]{Definition}
\theoremstyle{remark}
\newtheorem{remark}[theorem]{Remark}
\newtheorem{observation}[theorem]{Observation}

\title{Phase-Filtered Ramanujan Sums\\and the Spectral Gate}

\author{Alexander S.\ Petty}

\email{alexander.petty@gmail.com}
\date{June 2024}
\subjclass[2020]{11A63, 11B83, 42A16}

\begin{document}
\begin{abstract}
For primitive-root primes, the autocorrelation
$R(\ell)$ of the base repetend can be interpreted as a
phase-filtered Ramanujan sum, where the filter is the
cross-spectral function $G$ of the bin partition
evaluated along the multiplicative orbit of the base.

The continued fraction expansion of $p/b$ has a tail
determined by $r = p \bmod b$, partitioning primes into
$\varphi(b)$ spectral classes. Within each class, the
coarse bin pattern (the number of large bins,
$R = (p{-}1) \bmod b$) is determined by the class. The
continued fraction tail controls which lags produce
digit matches and which produce complete cancellation.
\end{abstract}
\maketitle

% ============================================================
\section{Introduction}

This paper expresses the autocorrelation as a phase-filtered Ramanujan sum.

Given a prime $p$ not dividing~$b$ with repetend
$(d_0, d_1, \ldots, d_{L-1})$ for $1/p$ in base~$b$, the
\emph{autocorrelation of the repetend} at lag~$\ell$ is
\[
R(\ell) = |\{j \in \{0, \ldots, L{-}1\} : d_j = d_{(j+\ell) \bmod L}\}|,
\]
the number of positions~$j$ at which the repetend digit equals the
digit $\ell$~positions later (cyclically).

The classical Ramanujan sum~\cite{ramanujan}
\[
c_q(n) = \sum_{\substack{k=1 \\ \gcd(k,q)=1}}^{q}
e^{2\pi i kn/q}
\]
assigns a uniform weight to each coprime residue. For a prime $q =
p$, $c_p(n) = -1$ for all $n \not\equiv 0 \pmod{p}$, and
$c_p(0) = p - 1$. The Ramanujan expansion of arithmetic functions~\cite{hardy_ramanujan}
connects these sums to the Riemann zeta function through the
identity
\[
\sum_{q=1}^{\infty} \frac{c_q(n)}{q^s}
= \frac{\sigma_{1-s}(n)}{\zeta(s)}.
\]

The autocorrelation $R(\ell)$ of the base repetend, derived in
the autocorrelation formula~\cite{paper9}, also sums over residues
modulo $p$. But unlike $c_p$, it
does not treat all nonzero lags equally. Some lags produce partial
constructive interference ($R(\ell) > 0$) while others produce
exact cancellation ($R(\ell) = 0$). The mechanism governing this
discrimination is the \emph{phase geometry} of the digit bins.

This paper identifies $R(\ell)$ as a phase-filtered Ramanujan sum,
characterizes the filter through the continued fraction expansion
of $p/b$, and classifies all primes into finitely many spectral
classes by their residue modulo~$b$.

% ============================================================
\section{The Phase Filter}

\begin{definition}
The \emph{phase filter} at lag $\ell$ is the ratio
\[
\Gamma(\ell) = \frac{|\Sigma_0(\ell)|}{M(\ell)},
\]
where $\Sigma_0(\ell) = \sum_k G(k, -kb^{-\ell} \bmod p)$ is the
cross-spectral sum from the autocorrelation formula~\cite{paper9}
and
$M(\ell) = \sum_k \sum_d |\widehat{\mathbf{1}}_{B_d}(k)|\,
|\widehat{\mathbf{1}}_{B_d}(-kb^{-\ell} \bmod p)|$
is the corresponding magnitude sum.
\end{definition}

The phase filter satisfies $0 \le \Gamma(\ell) \le 1$. When
$\Gamma = 1$ (all phases aligned), the cross-spectral sum achieves
its maximum and $R(\ell)$ is large. When $\Gamma = 0$ (complete
cancellation), $R(\ell) = 0$ regardless of the magnitudes.

\begin{observation}\label{prop:filter-values}
For a prime $p$ where $b$ is a primitive root (verified
computationally for all primes $p \le 10000$ in bases
$b \in \{3, 6, 7, 10, 12\}$):
\begin{enumerate}
\item $\Gamma(0) = 1$ (self-lag: perfect alignment).
\item For $\ell \ne 0$: $\Gamma(\ell) \in \{0\} \cup
      [c_{\min}, c_{\max}]$ where $c_{\min} > 0$. The filter is
      either exactly zero or bounded away from zero.
\item The set $\{\ell : \Gamma(\ell) > 0, \, \ell \ne 0\}$ has
      cardinality determined by $r = p \bmod b$.
\end{enumerate}
\end{observation}

The phase filter converts the uniform Ramanujan sum $c_p(\ell) =
-1$ into the selective autocorrelation $R(\ell)$: it ``gates''
which lags contribute to the spectral structure and which are
suppressed by destructive interference.

% ============================================================
\section{The Continued Fraction Control}

The bin pattern of the digit function $\delta(r) = \lfloor br/p
\rfloor$ is governed by the continued fraction expansion of $p/b$.

\begin{proposition}\label{prop:cf-structure}
The continued fraction~\cite{khinchin} of $p/b$ has the form
\[
\frac{p}{b} = [q,\; a_1, a_2, \ldots]
= \left[q,\; \operatorname{CF}\!\left(\frac{b}{r}\right)\right],
\]
where $q = \lfloor p/b \rfloor$ and $r = p \bmod b$. The tail
$\operatorname{CF}(b/r)$ depends only on $r$ and $b$, not on $p$.
\end{proposition}

\begin{proof}
The Euclidean algorithm: $p = qb + r$, so $p/b = q + r/b$ and the
continued fraction continues with $b/r$.
\end{proof}

Since $b$ is fixed, the tail of the CF is determined by $r$ alone.
In base~$10$, the coprime residues $r \in \{1, 3, 7, 9\}$ produce
four distinct CF tails:

\begin{table}[h]
\centering
\begin{tabular}{ccl}
\toprule
$r = p \bmod 10$ & CF tail & Bin pattern \\
\midrule
$1$ & $[10]$       & uniform ($R = 0$ large bins) \\
$3$ & $[3, 3]$     & $R = 2$ large bins \\
$7$ & $[1, 2, 3]$  & $R = 6$ large bins \\
$9$ & $[1, 9]$     & $R = 8$ large bins \\
\bottomrule
\end{tabular}
\caption{Spectral classes in base $10$. The CF tail determines the
bin pattern and hence the phase geometry of $G$.}
\end{table}

% ============================================================
\section{Spectral Classes}

\begin{definition}
Two primes $p_1, p_2$ (not dividing $b$) belong to the same
\emph{spectral class} if $p_1 \equiv p_2 \pmod{b}$.
\end{definition}

Within a spectral class, all primes share the same CF tail
and the same coarse bin pattern ($R = (p{-}1) \bmod b$
large bins). The number of constructive lags depends on
$p$ itself ($p - b - 1$ for primitive-root primes), not
on the spectral class alone.

\begin{theorem}\label{thm:spectral-classes}
In base $b$, there are $\varphi(b)$ spectral classes, where
$\varphi$ is Euler's totient function. In base $10$, there are
four: $r \in \{1, 3, 7, 9\}$.
\end{theorem}

The spectral class of $p$ determines:
\begin{itemize}
\item the number of large bins: $r$,
\item the CF tail: $\operatorname{CF}(b/r)$,
\item the phase filter $\Gamma(\ell)$ (up to the orbit-dependent
      evaluation),
\item and therefore the spectral richness (number of nonzero
      $R(\ell)$ values for $\ell > 0$).
\end{itemize}

% ============================================================
\section{Ramanujan Sums and the Autocorrelation}

The Ramanujan sum $c_p(\ell)$ and the autocorrelation $R(\ell)$
are both sums over $(\mathbb{Z}/p\mathbb{Z})^{*}$ involving
additive characters. Their relationship can be stated precisely.

\begin{proposition}\label{prop:ramanujan-R}
For a prime $p$ where $b$ is a primitive root:
\[
R(\ell) = \frac{1}{p} \sum_{k=0}^{p-1}
G(k, -kb^{-\ell} \bmod p),
\]
while
\[
c_p(\ell) = \sum_{k=1}^{p-1} e^{2\pi i k\ell/p}.
\]

Both sums run over additive characters of
$\mathbb{Z}/p\mathbb{Z}$. The Ramanujan sum weights each
character uniformly ($e^{2\pi ik\ell/p}$). The autocorrelation
weights each character by the cross-spectral function $G$,
evaluated at a point determined by the multiplicative orbit of
$b$.
\end{proposition}

The autocorrelation $R(\ell)$ is therefore a Ramanujan-type sum
with a non-uniform weight function. The weight function is $G$,
which encodes the digit-bin geometry. The Ramanujan sum is the
special case where $G$ is replaced by the trivial weight (all
bins of size~$1$, i.e., the digit-partitioning case).

\begin{corollary}
For digit-partitioning primes, $G(k,k') = \delta_{kk'} \cdot
(p{-}1)$ and $R(\ell) = (p{-}1)\delta_{\ell 0}$. The phase
filter passes only the self-lag, and the autocorrelation reduces
to the identity contribution of the Ramanujan sum.
\end{corollary}

% ============================================================
\section{Toward a Modified Zeta Identity}

The classical Ramanujan expansion connects $c_q(n)$ to $\zeta(s)$.
If $R(\ell)$ can be expressed as a modified Ramanujan sum with a
correction from $G$, a modified zeta identity may exist.

Define the \emph{structural Ramanujan sum} at prime $p$:
\[
c_p^{G}(\ell) = \frac{1}{p} \sum_{k=0}^{p-1}
G(k, -kb^{-\ell} \bmod p).
\]
For primitive-root primes, $R(\ell) = c_p^{G}(\ell)$ by the
autocorrelation formula~\cite{paper9}. The generating function
\[
L^{G}(s) = \sum_{p} \frac{c_p^{G}(\ell)}{p^s}
\]
(summing over primes at fixed $\ell$) is a Dirichlet series
encoding how the phase-filtered autocorrelation distributes across
primes.

Whether $L^{G}(s)$ has meromorphic continuation, a functional
equation, or a relationship to $\zeta(s)$ analogous to the
classical Ramanujan identity is an open question. The spectral
class structure (Theorem~\ref{thm:spectral-classes}) implies that
$L^{G}$ decomposes into contributions from $\varphi(b)$ residue
classes, suggesting a connection to Dirichlet $L$-functions
$L(s, \chi)$ for characters $\chi$ modulo~$b$~\cite{iwaniec}.

This direction is speculative. The series is defined only
for primitive-root primes, and its analytic properties are
not established.

% ============================================================
\section{Remarks}

\subsection*{The spectral gate}

The phase filter $\Gamma(\ell)$ acts as a gate: it admits certain
Ramanujan lags into the autocorrelation and blocks others. The
gate is controlled by the digit-bin geometry, which is controlled
by the CF expansion of $p/b$. The gate is therefore a structural
property of the interaction between the prime and the base, not a
property of either alone.

For DP primes, the gate is maximally restrictive: only the self-lag
passes. For non-DP primes, the gate opens selectively. The spectral
class determines which lags are admitted.

\subsection*{Connection to earlier work}

The recursive digit resolution framework~\cite{rdr} studies how
the interaction between a base grid and a prime grid creates
structure at different resolution scales. The continued fraction
expansion of $p/b$ encodes these scales. The partial quotients are
the resolution levels of the Moir\'{e} pattern between the two
grids. Large partial quotients correspond to near-alignment (flat
spectrum, DP-like behavior). Small partial quotients correspond to
misalignment (structured spectrum, rich autocorrelation).

The spectral gate is the formal mechanism by which this resolution
interaction governs the coherence of the fractional field.

% ============================================================
\begin{thebibliography}{9}

\bibitem{paper9}
A.~S.~Petty,
The autocorrelation formula,
research note, 2023.

\bibitem{rdr}
A.~S.~Petty,
Recursive digit resolution,
unpublished notes.

\bibitem{ramanujan}
S.~Ramanujan,
On certain trigonometrical sums and their applications in the
theory of numbers,
\emph{Trans.\ Cambridge Philos.\ Soc.} \textbf{22} (1918), 259--276.

\bibitem{hardy_ramanujan}
G.~H.~Hardy,
Note on Ramanujan's trigonometrical function $c_q(n)$ and certain
series of arithmetical functions,
\emph{Proc.\ Cambridge Philos.\ Soc.} \textbf{20} (1921), 263--271.

\bibitem{khinchin}
A.~Ya.~Khinchin,
\emph{Continued Fractions},
Dover, 1997.

\bibitem{iwaniec}
H.~Iwaniec and E.~Kowalski,
\emph{Analytic Number Theory},
Amer.\ Math.\ Soc., 2004.

\end{thebibliography}

\end{document}
